\documentclass[12pt,a4paper]{article}

% ============================================================
% PENGATURAN DASAR
% ============================================================
\usepackage[utf8]{inputenc}
\usepackage[T1]{fontenc}
\usepackage[a4paper,margin=2.5cm]{geometry}
\usepackage{titlesec}
\usepackage{enumitem}
\usepackage{hyperref}
\usepackage{xcolor}
\usepackage{fancyhdr}

\hypersetup{
    colorlinks=true,
    linkcolor=black,
    urlcolor=blue,
    citecolor=black
}

\pagestyle{fancy}
\fancyhf{}
\fancyhead[L]{Jurnal Belajar --- Algorithms for Decision Making}
\fancyhead[R]{Hari 1}
\fancyfoot[C]{\thepage}

\titleformat{\section}{\Large\bfseries}{\thesection}{1em}{}
\titleformat{\subsection}{\large\bfseries}{\thesubsection}{1em}{}
\titleformat{\subsubsection}{\normalsize\bfseries}{\thesubsubsection}{1em}{}

% Jarak sebelum/sesudah judul dirapatkan (before, after)
\titlespacing*{\section}{0pt}{10pt}{4pt}
\titlespacing*{\subsection}{0pt}{8pt}{2pt}
\titlespacing*{\subsubsection}{0pt}{6pt}{2pt}

% Rapatkan jarak antar paragraf dan daftar itemize
\setlength{\parskip}{2pt}
\setlist[itemize]{topsep=2pt, itemsep=2pt, parsep=0pt}

% ============================================================
% METADATA ENTRI HARI INI
% ============================================================
\newcommand{\harike}{1}
\newcommand{\tanggalentri}{25 Juli 2026}
\newcommand{\babyangdibahas}{Bab 1 --- Pendahuluan: 1.1 Decision Making s.d. awal 1.2.1 Aircraft Collision Avoidance}

\title{\textbf{Jurnal Belajar Harian} \\ \large Algorithms for Decision Making \\ \normalsize Kochenderfer, Wheeler, dan Wray (MIT Press, 2022)}
\author{}
\date{Hari ke-\harike\ --- \tanggalentri}

\begin{document}

\maketitle
\thispagestyle{fancy}

\noindent\textbf{Materi hari ini:} \babyangdibahas

\vspace{0.5cm}
\hrule
\vspace{0.5cm}

% ============================================================
% BAGIAN 1: PEMAHAMAN PRIBADI
% ============================================================
\section{Pemahaman Pribadi}

\subsection{Apa yang Dipelajari}
\begin{itemize}[leftmargin=1.2cm]
    \item Konsep \textit{agent} sebagai entitas (fisik maupun perangkat lunak) yang bertindak berdasarkan observasi atas lingkungannya, lewat siklus observasi-aksi yang berulang.
    \item Empat sumber ketidakpastian yang jadi kerangka besar seluruh buku: ketidakpastian hasil, model, keadaan, dan interaksi antaragen.
    \item Kerangka \textit{agent}-\textit{environment} tersebut mulai diterapkan pada kasus nyata pertama: sistem penghindaran tabrakan pesawat (\textit{aircraft collision avoidance}), yang memberi peringatan manuver kepada pilot berdasarkan data transponder.
\end{itemize}

\subsection{Insight}
\begin{itemize}[leftmargin=1.2cm]
    \item Empat kategori ketidakpastian yang dipaparkan (hasil, model, keadaan, interaksi) terasa seperti kerangka klasifikasi yang akan terus dipakai berulang di bab-bab berikutnya, bukan sekadar daftar contoh.
    \item Contoh tabrakan pesawat langsung memperlihatkan bahwa observasi yang tidak sempurna (radar bisa gagal mendeteksi) bukan detail teknis kecil, melainkan alasan utama mengapa pendekatan probabilistik dibutuhkan sejak awal.
\end{itemize}

\subsection{Pertanyaan Terbuka}
\begin{itemize}[leftmargin=1.2cm]
    \item Dari empat sumber ketidakpastian, mana yang paling dominan pada kasus penghindaran tabrakan pesawat, ketidakpastian hasil (respons pilot) atau ketidakpastian interaksi (perilaku pesawat lain)?
    \item Bagaimana batas kuantitatif antara ``peringatan terlalu dini'' dan ``peringatan tepat waktu'' sebenarnya ditentukan? Apakah ini akan dijelaskan lebih rinci di bab-bab berikutnya?
\end{itemize}

\vspace{0.5cm}
\hrule
\vspace{0.5cm}


\end{document}